% Source: sources/returns/S6_04_MAPS_MANIFOLDS/sections/manifolds.tex:4396–4701
\paragraph{Theorem D.2.22 (Proved: finite initial data for sphere outputs).}
Consider the standard section-line construction over any of the twenty-one
rational elliptic profiles with Mordell--Weil group $\mathbb Z$ in D.2.2.
Write $G$ for its chosen generator, $\langle G,G\rangle=1/d$,
$P=aG$, $Q=bG$, and $M=\mathcal O(P-O)$. Use a nonzero holomorphic
linearization with effective vertical divisor $D$, and the following local
fillings: a zero-free original quotient or a free affine logarithmic
replacement at each additive fibre; an original quotient or a standard
smooth Mumford filling at each multiplicative fibre; and, if desired,
free logarithmic replacements and rank-one Mumford fillings at regular
fibres. The latter occur at the regular zeros of the linearization.
These are the filling types specified in D.2; a new additive degeneration
at a zero of the linearization is not one of these local types.

If the completed manifold is a homotopy six-sphere, then $D$ is one
reduced component of one multiplicative fibre $I_n$. Translation by $Q$
acts transitively on its components and
\[
             \langle P,Q\rangle=\frac1n,
             \qquad ab=\frac dn.                                      \tag{D.FI1}
\]
After simultaneously changing the signs of $P,Q$, take $a,b>0$.
There are exactly $172$ marked triples (profile, ordered section pair,
zero position) satisfying the resulting integral component criterion.
They are listed completely in
\texttt{records/manifolds\_sphere\_natural.json}.
This is a finite list of possible initial data, not a claim that all
its filling choices are spheres.

The same necessary reduction holds in the larger construction with
$M=\mathcal O(P-O+V)$ for an arbitrary integral vertical divisor $V$.
Modulo simultaneous translation of the chosen zero component, there
are $226$ marked (profile, ordered section pair, zero-fibre) candidates;
all $226$ component equations have integral solutions $V$.
Their remaining solution freedom is exactly the $Q$-invariant vertical
divisor lattice. This optional enlargement is distinguished from the
natural section-line construction in the preceding paragraph.

\begin{proof}
We give the local Euler identity, the rational rank obstruction, and the
finite integral enumeration separately. No logarithmic multiplicity is
bounded in any step.

\emph{1. The local line and the determinant.}
Near a multiplicative fibre choose a Tate coordinate with $q=t^n$;
holomorphic unit factors do not affect any valuation below. If $P=t^p$
times a unit, choose $0\le p<n$. The product defining the Tate theta
function gives
\[
 \theta(qz;q)=-z^{-1}\theta(z;q),\qquad
 \sigma(z)=\frac{\theta(z/P;q)}{\theta(z;q)},\qquad
 \sigma(qz)=P\sigma(z).
\]
Thus $\sigma$ is the meromorphic section with divisor $P-O$ of the
line with factor of automorphy $P$. At a generic point of component
$i$, represented by $z=t^i$ times a unit, its valuation in the flat
frame is
\[
 f_i=\min(0,i-p)\quad(0\le i<n),\qquad f_{i+n}=f_i+p.
\]
Indeed, in this range only the factor $1-z/P$ of the numerator can
have negative valuation; the other theta factors are units. The
meromorphic section has no vertical divisor in $\mathcal O(P-O)$.
Consequently the flat frame, expressed in an extending local frame,
has valuation $F_i=-f_i$. Adding $V=\sum v_iC_i$ changes this to
$F_i=-f_i+v_i$, where $v_{i+n}=v_i$. Put $A=-p$, so
$F_{i+n}=F_i+A$.

If translation by $Q$ shifts the component index by $B$, let $w$ be
the valuation of its flat linearization multiplier. Comparing the
extending frames in $M$ and $t_Q^*M$ gives, with no rational
intersection-theory ambiguity,
\[
 z_i=\operatorname{ord}_{C_i}(D)=w+F_i-F_{i+B}.
\]
For a quasi-periodic sequence with increment $A$, telescoping yields
\[
 \sum_{i=0}^{n-1}(F_{i+B}-F_i)=AB,
 \qquad \sum_i z_i=nw-AB.                                    \tag{D.FI2}
\]
The identity holds for every integer $B$, by division by $n$ and
then telescoping the remaining finite interval. The actual valuation
vectors of the two long periods are $(n,-A)$ and $(B,-w)$: the latter
uses the inverse of the lift $M_{Qz}\to M_z$. Reversing the second
short-circle coordinate changes them to $(n,A),(B,w)$. Their index
is therefore $|nw-AB|$. For an effective nonzero local divisor this
equals $\sum_i z_i$. A standard unimodular triangular Mumford
cellulation has twice this many zero-dimensional torus strata, and
all its positive-dimensional strata have Euler number zero. Hence
\[
                  \chi(U_{I_n})=2\sum_i z_i.                  \tag{D.FI3}
\]
This proves the local equality directly from the extending line
frames, rather than identifying an arbitrary support function with
the line. For example, $n=2,p=B=1,w=0$ gives
$F=(1,0)$, $A=-1$, $(z_0,z_1)=(1,0)$, and index one, precisely
the component-supported zero of D.2.11.

\emph{2. Local rational facts used in the rank obstruction.}
Let $F=T^4$, $A_i$ its monodromy at a deleted point, $L_i$ the
mapping-torus boundary, and $U_i$ a filling. Put
\[
 f_i=\dim H_1(F;\mathbb Q)_{A_i},\quad
 g_i=\dim H_2(F;\mathbb Q)_{A_i},\quad h_i=b_1(U_i).
\]
For every allowed local type, nearby-fibre specialization generates
$H_1(U_i;\mathbb Q)$; a multiple meridian is a fibre loop. Thus
$h_i\le f_i$. Here are the actual local models and the degree-two
fact needed below.

A rank-one Mumford cap is topologically a cycle of nodal elliptic
components times $T^2$: the degree-zero ruled line has topologically
trivial Chern class, and its translation gluing is isotopic to the
identity. The surviving circle of the nodal curve, its fibre-area
class, and the two transverse circles show that
$H^2(U_i;\mathbb Q)\to H^2(F;\mathbb Q)^{A_i}$ is onto.
For a rank-two triangular cap, take the universal toric model with
rays $(v,1)$, $v\in\mathbb Z^2$, and the standard periodic triangular
fan. Its period lattice $L\subset\mathbb Z^2$ acts by
\[
 (z_1,z_2,t)\longmapsto
 (u_{\ell,1}(t)t^{\ell_1}z_1,
  u_{\ell,2}(t)t^{\ell_2}z_2,t),\qquad \ell\in L.
\]
Here $u:L\to\mathcal O^*(\Delta)^2$ is a character. It extends to
$\mathbb Z^2$: put the inclusion in Smith form and take holomorphic
roots of its two nowhere-zero unit values on the disc. The extended
period action uses the same fan and shifts its dual complex by
$\mathbb Z^2$. Its quotient is a minimal hexagonal cap. The original
cap covers this quotient with deck group $\mathbb Z^2/L$. This
action is free even on the central fibre, because its action on
$\mathbb R^2/L$ is a nontrivial translation. This is an ordinary
finite covering, not the toric power map discussed in C.1.
Pull back the four independent specialization classes of D.1.2. On the regular torus
the pullback is an isomorphism over $\mathbb Q$ onto the invariant
space, of dimension four; hence it again gives surjectivity.
Subdividing the toric cellulation preserves these pulled-back classes.
For free finite affine quotients, rational transfer and averaging
give the same assertion. This includes higher logarithmic orders,
the good-reduction caps, and the semistable covers of odd $I_n^*$.

There is a separate check for an original cap; it cannot be replaced
without checking its component Chern class. Its central model is
the circle bundle of $M$ over the mapping torus of the component
translation of the Kodaira curve. Since the linearization is a unit
there, its multidegree vector $c$ is invariant under that translation.
If $c\ne0$, the rational circle-bundle transgression kills the
$\delta$ circle, so $h_i=f_i-1$. If $c=0$, let $\rho$ be the number
of component orbits. On the invariant component subspace, the
Kodaira intersection matrix is negative semidefinite with kernel
the whole fibre and rank $\rho-1$. The map
$H_2(U_i)\to H_2(U_i,L_i)$ has this rank: pair its image with the
component four-cycles (the component times the two circle factors).
Thus the cokernel of $H_2(L_i)\to H_2(U_i)$ has dimension at least
$\rho-1$. If $c=0$ and $\rho=1$, the invariant fibre-area class
and the circle classes give the same degree-two surjectivity as
above. These statements also follow directly for a multiplicative
curve from its cycle of spheres times $T^2$; for an additive curve
use its tree of spheres. Finite component permutations are handled
by averaging, which is legitimate over $\mathbb Q$.

\emph{3. The global rank budget.}
Delete $N$ discs containing every singular fibre, every zero of the
linearization, and every extra logarithmic site, and write $E$ for
the resulting torus bundle. Its fibre coinvariants in degree one
have dimension $r=1$. To verify this without assuming a special
section pair, observe that $\mathcal O(aG-O)$ differs from
$\mathcal O(G-O)^{\otimes a}$ by a vertical line, trivial on this
punctured base. The elliptic-plus-$\delta$ monodromy block therefore
scales its last row by $a$. For $a\ne0$ this does not change its
rational image rank. The exact twenty-one canonical blocks saved
by the program all have rank three. All full monodromy differences
lie in these first three coordinates, and hence have precisely
this image. Vertical modifications and the $Q$ column cannot
change the conclusion.

Also $v_2=\dim H_2(F;\mathbb Q)_{\pi_1}\ge1$. Explicitly,
in the marking $(e_1,e_2,\delta,\gamma)$ the actual monodromy
has the following blocks, where $z$ is the section cocycle of $G$:
\[
 A=\begin{pmatrix}
 B&0&bz\\ -az^tJB&1&c\\0&0&1
 \end{pmatrix},\qquad
 J=\begin{pmatrix}0&1\\-1&0\end{pmatrix},\qquad B\in SL_2(\mathbb Z).
\]
The same vertical-line trivialization gives the row multiplied by
$a$, while the Abel--Jacobi lift of $Q=bG$ gives the column $bz$.
The central entry $c$ is unrestricted. Direct block multiplication
shows that $A^t\Omega A=\Omega$, where
\[
 \Omega=\begin{pmatrix}aJ&0\\
 0&\begin{matrix}0&-b\\ b&0\end{matrix}\end{pmatrix}.
\]
In the only potentially nonzero cross term,
$B^tJz=(-z^tJB)^t$, so the two terms cancel. Thus
$\Omega=a e_1^*e_2^*+b\gamma^*\delta^*$ is a nonzero invariant,
independently of $c$ and of any vertical modification. The program
also verifies this identity symbolically. The case $a=0$ or $b=0$ cannot have an effective
nonempty $D$: norm over a component-preserving translation makes
such a divisor a base divisor of degree $\langle P,Q\rangle=0$.
It therefore cannot give Euler number two.

The graph-base homology sequence gives
\[
 b_1(E)=N,\qquad b_2(E)=v_2+4N-7,\qquad b_2(L_i)=g_i+f_i.
\]
If $H_1(X;\mathbb Q)=H_2(X;\mathbb Q)=0$, the first MV map is
onto and its kernel has dimension $\sum_i(f_i-h_i)$. This kernel
is a quotient of $H_2(X;\mathbb Q)$, so every $h_i=f_i$.
In particular every original cap has $c=0$. The degree-two map
from the boundary is also onto each cap; the component intersection
calculation then forces $\rho=1$ for every original cap.
All remaining caps consequently satisfy the degree-two invariant
surjectivity proved in step 2, and so $b_2(U_i)\ge g_i$.
After first mapping onto $\bigoplus H_2(U_i)$, the boundary kernel
has dimension at most $\sum f_i$. It must still surject onto
$H_2(E)$, since $H_2(X)=0$. Therefore
\[
 v_2+4N-7\le\sum_i f_i,
 \qquad\sum_i\operatorname{rank}(A_i-I)\le7-v_2\le6.          \tag{D.FI4}
\]

Let $F_a,F_m$ count additive and multiplicative fibres of the
rational elliptic surface. The Shioda--Tate formula, $\rho(S)=10$,
and Mordell--Weil rank one give root rank seven. Comparing this
with $\chi(S)=12$ gives
\[
                            2F_a+F_m=5.
\]
This is also checked independently for every profile in the program;
see \cite[Theorem 6.3 and Corollary 6.13]{manSS10}.
Every additive monodromy difference has rank at least two, and
every multiplicative one rank at least one. A rank-two Mumford
zero at a multiplicative fibre raises its rank to two; a zero at
a regular fibre contributes an additional rank one. All original
and logarithmic caps have Euler number zero, as do rank-one caps.
Thus $\chi(X)=2$ and (D.FI3) force at least one rank-two zero,
already consuming the sole extra rank allowed by (D.FI4).
There are no regular zeros and no second multiplicative zero.
Equation (D.FI3) now says the effective divisor at that fibre has
total coefficient one. It is one reduced component, as asserted.

\emph{4. Norm, height, and finite enumeration.}
Choose an integer $R$ such that $RQ$ meets the identity component
of every reducible fibre. Norm the line identity over the first
$R$ iterates of $t_Q$. The resulting divisor has zero component
multidegrees, and is therefore a sum of whole fibres. Since it is
effective and supported on the orbit of the single component,
$Q$ must act transitively on this $I_n$. Its norm is $(R/n)F$.
The normed line is
$M\otimes t_{RQ}^*M^{-1}$; the vertical correction cancels because
$RQ$ fixes every component. The height formula with a narrow second
section gives the base degree
$\langle P,RQ\rangle=R\langle P,Q\rangle$; all its local correction
terms vanish. Hence (D.FI1). The divisor and height comparison is
integral: a vertical divisor of zero multidegree is a whole-fibre
multiple, by the corank-one component intersection matrix, and
$\operatorname{Pic}^0(S)=0$.

Here is the entire finite algorithm, with no implicit parameter
bound. For each multiplicative $I_n$ with $n\mid d$, list the
positive ordered divisors $ab=d/n$. Require that the component
of $Q$ generate $\mathbb Z/n$. For each Kodaira fibre let $I$ be
its integral component intersection matrix, $T$ the component
pullback permutation of $Q$, and $c_P=e_{P\text{-component}}-e_0$.
For a candidate divisor vector $z$, solve
\[
              (1-T)I V=Iz-(1-T)c_P.                         \tag{D.FI5}
\]
At the chosen zero fibre, $z$ is one component; elsewhere $z=0$.
For the natural construction require $V=0$ and test every possible
zero component. For the larger construction solve (D.FI5) by
unimodular Smith operations. Since $I$ is invertible modulo the
invariant whole fibre, its solution kernel is exactly $\ker(1-T)$.
No additional numerical condition is omitted: (D.FI5) makes
$J\otimes\mathcal O(-D)$ a base pullback, and its norm has degree
zero by (D.FI1), so the base pullback is trivial. Thus every retained
row actually has this effective linearization divisor.

The program checks the affine component diagrams, their primitive
multiplicity kernels and the translation permutations, then solves
every (D.FI5) integrally. The following counts retain marked fibre
positions; the first number uses $V=0$, the second allows $V$:
\[
\begin{array}{c|rr@{\qquad}c|rr}
43&6&6&43/II&2&2\\
45&8&17&45/II&4&9\\
46&6&9&46/II&2&3\\
47&14&16&47/II&6&8\\
47/III&8&10&49&10&10\\
49/II&2&2&49/III&4&4\\
50&10&15&50/II&2&3\\
50/IV&4&6&55&14&23\\
55/II&6&11&56&26&28\\
56/II&10&12&56/III&14&16\\
56/IV&14&16&\text{total}&172&226
\end{array}
\]
For $n>1$, the natural component equation in fact forces
$P$ and $Q$ to meet the same adjacent component, labelled $1$ or
$-1$, and the zero to be $C_0$; the code tests this directly,
including the doubled edges for $I_2$. The broader count normalizes
the zero component to $C_0$ by a power of $Q$; translating $M$
changes it only by an allowed vertical divisor. This explains
the different marking convention in the two counts.
\end{proof}
